\documentclass[11pt,a4paper]{article}

% ===== XeLaTeX + Korean =====
\usepackage{kotex}
\usepackage{fontspec}
\setmainfont{Times New Roman}
\setmainhangulfont{AppleMyungjo}
\setsanshangulfont{Apple SD Gothic Neo}

% ===== Layout =====
\usepackage[a4paper,top=18mm,bottom=18mm,left=20mm,right=20mm]{geometry}
\usepackage{fancyhdr}
\setlength{\headheight}{24pt}
\usepackage{enumitem}
\usepackage{mdframed}
\usepackage{array}
\usepackage{booktabs}

\setlength{\parindent}{1em}
\linespread{1.18}

% ===== Header / Footer =====
\pagestyle{fancy}
\fancyhf{}
\renewcommand{\headrulewidth}{0.6pt}
\fancyhead[L]{\sffamily\bfseries 2026 고1 영어 변형 — 지문 24 정답 및 해설}
\fancyhead[R]{\sffamily 도파민과 SNS}
\renewcommand{\footrulewidth}{0pt}
\fancyfoot[C]{\framebox[16mm][c]{\thepage}}

% ===== helpers =====
\newcommand{\qno}[1]{\par\vspace{10pt}\noindent{\sffamily\bfseries #1.}\ }
\newcommand{\instr}[1]{{\sffamily #1}\par\vspace{3pt}}
\newmdenv[linewidth=0.6pt,roundcorner=3pt,innertopmargin=7pt,innerbottommargin=7pt,
          innerleftmargin=9pt,innerrightmargin=9pt,skipabove=5pt,skipbelow=5pt]{pbox}

\begin{document}

\begin{center}
{\fontsize{15}{17}\selectfont\sffamily\bfseries [지문 24] 정답 및 해설 — 도파민과 SNS}
\end{center}
\vspace{6pt}

% ===== 정답 표 =====
\begin{center}
\renewcommand{\arraystretch}{1.4}
\begin{tabular}{cccc}
\toprule
문항 & 유형 & 정답 & 배점 \\
\midrule
1 & 빈칸추론 & ② & 2점 \\
2 & 어휘 (문맥상 부적절) & ③ & 2점 \\
3 & 서술형 — 해석 & 모범답안 참조 & 4점 \\
4 & 영어 제목 & ⑤ & 2점 \\
\bottomrule
\end{tabular}
\end{center}

\vspace{8pt}

% ===================== Q1 해설 =====================
\qno{1}\instr{빈칸추론 — 정답: ②\, equilibrium}

\noindent\textbf{정답 근거:} 빈칸 직후 ``or balance''라는 동격 표현이 등장하여 빈칸이 `균형'을 의미하는 단어임을 직접 알려 준다. 또한 지문 앞부분에서 뇌가 항상 `homeostasis(항상성)'를 추구하며 이것이 ``simply means balance''라고 설명하므로, 빈칸에는 균형·평형을 뜻하는 \textbf{equilibrium}이 가장 적절하다.

\vspace{4pt}\noindent\textbf{오답 소거:}
\begin{itemize}[leftmargin=2em,itemsep=1pt,topsep=2pt]
  \item ① stimulation (자극) — 도파민 상승의 원인이지, 뇌가 달성하려는 상태가 아니다.
  \item ③ acceleration (가속) — 지문의 논지(균형 회복)와 반대 방향이다.
  \item ④ gratification (만족) — 도파민 상승 직후의 감각이나, 빈칸 뒤 ``or balance''와 의미가 맞지 않는다.
  \item ⑤ withdrawal (금단/철수) — 도파민 감소 이후의 부작용 상태이며, 뇌가 능동적으로 달성하려는 상태가 아니다.
\end{itemize}

% ===================== Q2 해설 =====================
\qno{2}\instr{어휘 — 정답: ③\, avoiding}

\noindent\textbf{정답 근거:} 지문은 뇌와 몸이 항상 `homeostasis(항상성)', 즉 균형을 \textbf{추구(seek)}한다고 설명한다. 따라서 ③의 ``avoiding(회피하다)''은 문맥과 정반대이며, \textbf{seeking}으로 바꾸어야 올바르다.

\vspace{4pt}\noindent\textbf{나머지 어휘 확인:}
\begin{itemize}[leftmargin=2em,itemsep=1pt,topsep=2pt]
  \item ① occurring — ``뇌에서 일어나고 있는'' 것을 묘사; 문맥상 적절.
  \item ② incredibly — 도파민이 ``믿을 수 없이'' 빠르게 상승한다는 의미; 적절.
  \item ④ baseline — 기준 수준 아래로 도파민이 떨어진다는 의미; 적절.
  \item ⑤ rebalance — 균형을 다시 맞추기 위해 도파민이 하락한다는 의미; 적절.
\end{itemize}

% ===================== Q3 해설 =====================
\qno{3}\instr{[서술형] 해석 — 모범답안 및 채점 포인트}

\noindent\textbf{대상 문장:}\\
\hspace{1em}\textit{the dopamine then has to quickly drop an equal amount below your baseline level in order to rebalance, making you feel even worse than before you started scrolling.}

\vspace{6pt}\noindent\textbf{모범답안:}\\
그러면 도파민은 재균형을 이루기 위해 당신의 기준 수준보다 동일한 양만큼 빠르게 떨어져야 하며, 이는 스크롤을 시작하기 전보다 더 나쁜 기분이 들게 만든다.

\vspace{6pt}\noindent\textbf{채점 포인트 (각 1점, 총 4점):}
\begin{enumerate}[leftmargin=2em,itemsep=2pt,topsep=2pt]
  \item ``drop an equal amount below your baseline level'' → 기준 수준보다 동일한 양만큼 하락/떨어지다 (기준 수준·동일한 양 모두 포함 시 인정)
  \item ``in order to rebalance'' → 재균형을 이루기/위해 (목적의 의미 전달 시 인정)
  \item ``making you feel even worse'' → 더 나쁜 기분이 들게 만들다 (비교급 ``even worse'' 반영 시 인정)
  \item ``than before you started scrolling'' → 스크롤을 시작하기 전보다 (비교 대상 명시 시 인정)
\end{enumerate}

\vspace{8pt}
\noindent{\sffamily\bfseries 4. 제목 - 정답: ⑤}
\par 글은 SNS 스크롤로 도파민이 빠르게 상승한 뒤, 항상성을 회복하기 위해 기준치 아래로 급락하면서 더 나쁜 기분이 생긴다는 과정을 설명한다. 따라서 ⑤가 가장 적절하다.

\end{document}
